\part{Canon, Sony et Nikon : principes communs, variantes réelles}
\section{Comment lire les informations propres à une marque}
La marque n'est pas une spécification de codec. Un modèle,
un firmware, un mode de rafale, un type d'obturateur, une
taille RAW et un réglage de compression peuvent changer
les données. Il faut constituer une matrice de support
\emph{modèle $\times$ mode $\times$ firmware}, pas une
simple liste d'extensions.

Une documentation de rétro-ingénierie est souvent la
meilleure explication publique de détails binaires,
mais elle doit être corroborée par des fichiers et des
implémentations. Une ancienne constante noire ou un
offset de MakerNote peut être valide sur un modèle et
faux sur son successeur. Un nouveau boîtier peut réutiliser
le conteneur tout en introduisant une nouvelle calibration.

\section{Canon : CR2, CR3, C-RAW et styles de rendu}
\subsection{CR2 : une organisation héritée de TIFF}
Les travaux de Laurent Clévy décrivent le CR2 comme une
structure TIFF enrichie, avec répertoires, métadonnées,
aperçus et données RAW. Ils expliquent aussi le décodage
JPEG prédictif sans perte et le réassemblage de tranches
rencontrés dans les fichiers étudiés~\cite{clevycr2}.
Ce JPEG sans perte n'est pas le JPEG à DCT avec pertes
d'une photographie livrée au client.

Il faut donc séparer trois opérations : trouver la charge
utile, décompresser ses échantillons, puis remettre les
échantillons dans leur géométrie capteur. Un tableau de
la bonne taille peut malgré tout avoir ses segments
permutés. Les modes réduits ou particuliers ne doivent
pas être traités comme une mosaïque pleine taille sans
vérifier leur organisation.

\subsection{CR3 : conteneur à boîtes et codec CRX}
La description publique CR3 de Clévy identifie notamment
\texttt{ftyp}, \texttt{moov}, des pistes, des données
\texttt{mdat}, des aperçus et des blocs de métadonnées
Canon. Les blocs CMT peuvent embarquer des structures
de type TIFF. La famille CRX comprend des variantes
prédictives et des variantes à ondelettes et quantification
pour C-RAW ; les détails observés sont documentés comme
rétro-ingénierie, non comme une norme officielle Canon~\cite{clevycr3}.

Le décodeur CRX de LibRaw contient lui-même des structures
de tuiles, plans, sous-bandes et paramètres de quantification~\cite{librawcrx}.
Le lecteur doit donc interpréter les descripteurs pour
savoir quel procédé appliquer. Déduire le mode du seul
suffixe CR3 est insuffisant. Le CR3 n'est pas non plus un
JPEG renommé : il peut seulement en contenir des aperçus.

\subsection{RAW contre C-RAW : ce que dit Canon et ce qu'on infère}
Canon présente C-RAW comme une variante de taille réduite
conservant la souplesse de développement RAW, et signale
que certaines remontées importantes de sous-exposition
peuvent révéler davantage de bruit. L'article constructeur
ne publie pas une description exhaustive de l'algorithme
binaire~\cite{canonraw}.

Les analyses et le code CRX apportent une description
plus technique de certains modes. Il ne faut pas convertir
un comportement observé sur une génération en garantie
pour toutes les futures versions. Pour tester un débruiteur
ou des ombres extrêmes, consigner explicitement RAW ou
C-RAW et éviter de mélanger les deux sans contrôle.

\subsection{Picture Style, Camera Neutral et Adobe Standard}
Canon distingue des styles de rendu tels que Standard,
Neutral et Faithful et décrit leurs objectifs visuels.
Le manuel du R6 Mark III distingue aussi force, finesse
et seuil de netteté~\cite{canonstyle}. Ces paramètres
participent à un rendu ; ils ne sont pas les coefficients
d'une simple matrice RGB à copier dans un décodeur.

Le style Neutre du boîtier, un DCP Adobe Camera Neutral
et un profil Adobe Standard sont trois objets différents.
Un profil Camera Matching cherche un rapprochement de
rendu, pas nécessairement une reproduction bit à bit du
moteur constructeur. Les noms n'impliquent pas non plus
une neutralité physique parfaite sous tous les éclairages.

Le JPEG intégré au RAW peut refléter le style choisi sur
le boîtier alors que le logiciel tiers développe la
mosaïque avec son propre profil par défaut. L'utilisateur
voit alors une transition entre l'aperçu Canon et le
rendu tiers, sans que les mesures RAW aient été changées.

\subsection{Exemple local IPR : calibrations à ne pas universaliser}
L'audit des six scènes Canon du projet a notamment
examiné une nouvelle version de ColorData, des zones
optical-black et la précision des faibles valeurs en
sortie. Une erreur de sélection de bordure peut biaiser
le noir ; une mauvaise interprétation du blanc change
la normalisation. La comparaison à Lightroom a aussi
isolé des différences de tables et de courbe.

Ces résultats datés sont dans
\href{../validation/analyses/analyse-base-colorimetrique.md}{le rapport évolutif de base colorimétrique}.
Les liens PDF relatifs dépendent du lecteur ; le chemin
depuis la racine est
\path{docs/validation/analyses/analyse-base-colorimetrique.md}.
Le cours retient la méthode, pas des valeurs à appliquer
à tous les RAW Canon. Aucun nouveau correctif n'est
introduit par la rédaction de ce document.

\section{Sony : ARW, variantes de compression et prudence spectrale}
\subsection{Ce qui est déjà documentable}
Les fichiers ARW relèvent d'une famille de structures
et de codecs Sony, avec métadonnées et aperçus.
Le code de RawSpeed fournit des exemples de décodage
de certaines variantes ARW, mais le nom d'une classe
ne démontre pas le support de tous les modes Sony~\cite{rawspeed,sonyarw2}.
Un lecteur doit reconnaître le modèle et la variante
avant d'interpréter profondeur, emballage et compression.

\subsection{Modes du Sony Alpha 7 IV : un exemple précis}
Le manuel Sony du 7 IV décrit des choix non compressé,
compressé sans perte L/M/S et compressé. Il précise que
M et S réduisent la résolution par rapport à L, et
mentionne une condition de firmware pour ces choix~\cite{sonymodes}.
Il ne faut pas lire « compression sans perte M » comme
« toutes les mesures de pleine résolution sont conservées » :
la compression peut être sans perte sur une représentation
déjà réduite spatialement.

Cette distinction s'applique bien au-delà de Sony :
la garantie d'un codec porte sur son entrée, pas sur
toutes les opérations antérieures de la chaîne caméra.
Pour un protocole rigoureux, choisir la pleine résolution
et documenter le mode, sans supposer qu'un autre modèle
dispose des mêmes menus.

\subsection{Un exemple ARW2 de codage par petits groupes}
Dans la variante illustrée par le décodeur SonyArw2
de RawSpeed, un groupe de seize échantillons consomme
128 bits, avec extrema, positions et valeurs résiduelles
compactes. Le code impose des invariants et reconstruit
les valeurs selon la plage du groupe~\cite{sonyarw2}.
Ce mécanisme est un exemple concret de codage adaptatif
à la dynamique locale, pas une description de tous les ARW.

Il explique pourquoi des erreurs ou compromis peuvent
être corrélés à des groupes et à de forts contrastes.
Les essais doivent inclure arêtes lumineuses, ombres
remontées et textures fines, et pas seulement des
aplats correctement exposés. Aucune estimation de
qualité Sony globale n'est déduite ici d'un seul codec.

\subsection{Couleur et futurs tests Sony dans IPR}
Une matrice Canon, ses gains et son profil ne peuvent
pas être réutilisés parce qu'un Sony aurait lui aussi
un CFA Bayer. Les réponses spectrales, calibrations,
normalisations et rendus diffèrent. Le futur corpus
Sony doit contenir un RAW plein format non réduit,
son mode exact, un rendu neutre de référence clairement
nommé et, idéalement, une charte sous éclairages connus.

\paragraph{À compléter avec des fichiers réels.} Inventorier
modèle/firmware, modes compressés et sans perte, géométrie
native, profondeur utile, niveaux, orientation, gains,
profil et effet des styles Sony. Cette section donne
une base documentée ; elle ne revendique aucune validation
de fichiers Sony dans la campagne Canon actuelle.

\section{Nikon : NEF et coexistence de procédés}
\subsection{NEF ne décrit pas un codec unique}
Les NEF utilisent des organisations de la famille TIFF
avec données Nikon et variantes de compression. Le
décodeur Nikon de RawSpeed illustre notamment des
tables, prédictions et décodages de différences pour
les variantes qu'il traite~\cite{nikondecoder}.
Ces mécanismes ne doivent pas être appliqués à tout
fichier Nikon sur le seul nom de marque.

\subsection{Exemple Z6/Z7 : profondeur et réversibilité}
Le manuel Nikon Z6/Z7 distingue compression sans
perte, compression non réversible et non compressé,
ainsi que des choix 12 ou 14 bits~\cite{nikonmodes}.
Les deux axes sont distincts : un RAW 14 bits avec
pertes et un RAW 12 bits sans perte ne se comparent
pas avec une seule étiquette « qualité ».

Le nombre de bits indique une capacité de codage,
pas à lui seul la dynamique effective du capteur.
Un choix de mode peut aussi modifier vitesse et
compatibilité ; il faut consulter le manuel du
modèle et son firmware au lieu d'extrapoler celui
d'un Z6 ou d'un Z7 à un Z9.

\subsection{Exemple Z6 III : haute efficacité}
Le manuel du Z6 III distingue compression sans
perte et deux variantes haute efficacité, avec
des tailles et objectifs de qualité différents~\cite{nikonefficiency}.
Le nom marketing « haute efficacité » ne suffit
pas à identifier une structure binaire ni à
garantir la réversibilité. La prise en charge
doit être vérifiée pour la bibliothèque et sa
version réellement utilisées.

Une bibliothèque peut ouvrir les métadonnées
et l'aperçu sans savoir développer un mode
de compression nouveau. Le test pertinent
demande les pixels natifs et leur cohérence,
pas seulement une icône visible dans l'explorateur.

\subsection{Picture Control et futurs tests Nikon}
Les réglages de rendu Nikon, comme les Picture
Controls, doivent être distingués de la mosaïque,
des métadonnées et des profils tiers. Pour une
comparaison, préciser si la référence est un
JPEG boîtier, un développement NX Studio ou
un développement Lightroom avec profil défini.

\paragraph{À compléter avec des fichiers réels.} Tester
au moins les modes pertinents du modèle retenu,
la profondeur, les dimensions, les aperçus,
l'orientation, les gains et la calibration.
Les formats de certaines gammes, comme NRW,
doivent faire l'objet d'une vérification séparée ;
ils ne sont pas couverts par une recette NEF
automatique supposée universelle.

\section{Synthèse inter-marques et limites de connaissance}
\begin{longtable}{p{3.2cm}p{4cm}p{7.5cm}}
\toprule Question & Invariant utile & Variable à vérifier \\ \midrule
Qu'est-ce qui est mesuré ? & Signal du capteur & CFA, profondeur, prétraitements et lecture. \\
Comment le fichier s'organise ? & Index et charges utiles & TIFF, boîtes, versions et offsets. \\
Que fait le codec ? & Reconstruit une représentation & Prédiction, transformée, pertes et mode réduit. \\
Comment retrouver les couleurs ? & Calibration nécessaire & Matrices, gains, tables et illuminants. \\
Quel rendu veut-on ? & Choix explicite & Style boîtier, profil tiers ou cible physique. \\
Peut-on faire confiance au décodage ? & Tests reproductibles & Modèle, firmware, mode et version bibliothèque. \\
\bottomrule
\end{longtable}
Le bon investissement logiciel consiste à généraliser
les contrats et les contrôles, pas à généraliser des
constantes issues d'un seul boîtier. Les informations
manquantes doivent rester identifiées comme telles.
